\section{Lightweight Iterative Optimization Method}
\label{sec:liom}
With the flight horizon fixed, LIOM reconciles local collision-avoidance regions and motion limits by refining both the trajectories and the corridor geometry. External penalties~\cite{li2022liom} retain intermediate candidates for successive corridor constructions. The process has three stages: a relaxed QP supplies a candidate, numerical recovery qualifies it for control reconstruction, and complete-interval checks determine acceptance of the resulting motion.

\subsection{Two-Penalty Recovery Formulation}
For the Cartesian motion model, the kinematic defect is affine in the position and velocity variables. The two quantities that measure constraint recovery are
\begin{equation}
 D=\sum_{i=1}^N\sum_{k=0}^{K-1}\norm{\mathbf E_{i,k}}_2^2,\qquad
 S=\sum_{(i,j,k)\in\mathcal A}\delta_{ij,k}.
 \label{eq:recovery}
\end{equation}
At iteration $\nu$, the augmented objective is
\begin{equation}
 J^{(\nu)}=C+\frac{\rho_q^{(\nu)}}hD+\rho_s^{(\nu)}S.
 \label{eq:penalty}
\end{equation}
The quantities $D$ and $S$ have units of $\mathrm{m}^2$ and $\mathrm{m}$, respectively. The factor $1/h$ scales the kinematic penalty with the time discretization. The weights $\rho_q^{(\nu)}$ and $\rho_s^{(\nu)}$ regulate the recovery of kinematics and relative-corridor compatibility, respectively.

Let $\mathcal I=\{1,\ldots,N\}$, $\mathcal K=\{0,\ldots,K\}$ and $\mathcal K^-=\{0,\ldots,K-1\}$. Corridor bounds and $\mathcal A^{(\nu)}$ are fixed during an inner solve. The complete subproblem is
\begin{adjustwidth}{-\extralength}{0cm}
\begin{subequations}\label{eq:subproblem}\begin{align}
 \min_{\mathbf P,\mathbf V,\boldsymbol\delta}\quad&
 C(\mathbf V)
 +\frac{\rho_q^{(\nu)}}h\sum_{i\in\mathcal I}\sum_{k\in\mathcal K^-}\norm{\mathbf E_{i,k}}_2^2
 +\rho_s^{(\nu)}\sum_{(i,j,k)\in\mathcal A^{(\nu)}}\delta_{ij,k}
 \label{eq:nlpobjective}\\
 \mathrm{s.t.}\quad&
 \underline{\mathbf w}+r\one\leq\bp_{i,k}\leq\overline{\mathbf w}-r\one,
 && i\in\mathcal I,\ k\in\mathcal K,\label{eq:nlpworkspace}\\
 &\underline{\boldsymbol\gamma}_{i,k}\leq\bp_{i,k}\leq\overline{\boldsymbol\gamma}_{i,k},
 &&i\in\mathcal I,\ k\in\mathcal K,\label{eq:nlpstc}\\
 &-v_c\one\leq\bv_{i,k}\leq v_c\one,
 &&i\in\mathcal I,\ k\in\mathcal K,\label{eq:nlpvelocity}\\
 &-a_c\one\leq(\bv_{i,k+1}-\bv_{i,k})/h\leq a_c\one,
 &&i\in\mathcal I,\ k\in\mathcal K^-,\label{eq:nlpacceleration}\\
 &\bp_{i,0}=\mathbf s_i,\quad\bp_{i,K}=\mathbf g_i,\quad
 \bv_{i,0}=\bv_{i,K}=\mathbf0,
 &&i\in\mathcal I,\label{eq:nlpboundary}\\
 &\underline{\boldsymbol\eta}_{ij,k}-\delta_{ij,k}\one
 \leq\bp_{i,k}-\bp_{j,k},
 &&(i,j,k)\in\mathcal A^{(\nu)},\label{eq:nlprellower}\\
 &\bp_{i,k}-\bp_{j,k}\leq\overline{\boldsymbol\eta}_{ij,k}+\delta_{ij,k}\one,
 &&(i,j,k)\in\mathcal A^{(\nu)},\label{eq:nlprelupper}\\
 &\delta_{ij,k}\geq0,
 &&(i,j,k)\in\mathcal A^{(\nu)}.\label{eq:nlpslack}
\end{align}\end{subequations}
\end{adjustwidth}
In the objective, $i$ ranges over $\mathcal I$. Acceleration and $\mathbf E_{i,k}$ are substituted expressions. Workspace and STC bounds are intersected, and endpoint states are imposed by equal lower and upper variable bounds. Every constraint is linear or a variable bound. Since $C,D$ are convex quadratic and $S$ is linear, the fixed-corridor program is a convex quadratic program (QP). The affine double-integrator equations and hard relative boxes also retain this convex structure. The two relaxations preserve availability of intermediate feasible points when those hard requirements conflict; corridor reconstruction changes the local geometric region between solves.

\begin{Proposition}[Nonemptiness of the relaxed fixed-corridor subproblem]
\label{prop:nonempty}
Suppose every spatial box is nonempty, contained in the admissible workspace, and contains the prescribed start or goal at the corresponding endpoint. For the rest-to-rest task, the relaxed subproblem~\eqref{eq:subproblem} has a feasible point with finite objective.
\end{Proposition}
\begin{proof}
Choose any position inside each box, preserve the fixed endpoint positions, and set every velocity to zero. The velocity, acceleration, and rest-to-rest boundary bounds are satisfied. For each active triple, choose
\begin{equation}
 \delta_{ij,k}=\max\!\left\{0,\,
 \max_d(\underline{\eta}_{ij,k,d}-r_{ij,k,d}),\,
 \max_d(r_{ij,k,d}-\overline{\eta}_{ij,k,d})\right\}.
 \label{eq:feasiblepoint}
\end{equation}
Both relative inequalities are satisfied. All chosen states and slacks are finite, hence so is the augmented objective.
\end{proof}
Under the stated box conditions, Proposition~\ref{prop:nonempty} ensures that the relaxed local problem remains available for refinement. Numerical recovery and final trajectory acceptance are assessed separately by the procedures below. If construction or numerical solution fails, or the budget expires before acceptance, the attempt terminates with planning failure.

\subsection{Bounded Continuation and Corridor Reconstruction}
Bounded external-penalty continuation increases the cost of residual inconsistency while allowing the corridor geometry to adapt. Numerical recovery is reached when
\begin{equation}
 D\leq\varepsilon_q,\qquad S\leq\varepsilon_s.
 \label{eq:tolerances}
\end{equation}
Here $\varepsilon_q$ and $\varepsilon_s$ are aggregate defect and slack tolerances, with units $\mathrm{m}^2$ and $\mathrm{m}$, specified for the chosen grid and team scale. Starting from $\rho_q^{(0)}$ and $\rho_s^{(0)}$, each coefficient is updated independently:
\begin{equation}
 \rho_a^{(\nu+1)}=
 \begin{cases}
 \min\{\beta_a\rho_a^{(\nu)},\overline\rho_a\},&Q_a^{(\nu)}>\varepsilon_a,\\
 \rho_a^{(\nu)},&Q_a^{(\nu)}\leq\varepsilon_a,
 \end{cases}
 \quad (Q_q,Q_s)=(D,S).
 \label{eq:weights}
\end{equation}
The growth factors $\beta_a>1$ provide gradual continuation, and the upper bounds $\overline\rho_a$ limit the penalty scale. A finite unsuccessful candidate supplies the next reference and warm start. The initial guides remain available for static seed relocation.

A candidate satisfying~\eqref{eq:tolerances} and the hard constraints within tolerance $\varepsilon_h$ enters the control-integration step below. The resulting motion must satisfy the prescribed physical limits and clearance conditions before it is accepted. Otherwise, the candidate supplies the next reference. Iteration ends at the first accepted trajectory or when the prescribed iteration and computation budgets are exhausted. The acceleration-effort term regularizes these local solves; recovery and acceptance govern outer termination, and the final control correction can change its value.

\subsection{Exact Rest-to-Rest Output from Recovered Controls}
\label{sec:arcs}
Finite kinematic defects are handled directly at the control level. For one spatial coordinate of UAV $i$, collect candidate accelerations in $\mathbf a_i\in\R^K$. In the following two equations, $s_i,g_i,v_i^0,v_i^f$ denote components along this same coordinate. Constant-acceleration integration expresses terminal velocity and displacement as two control moments:
\begin{equation}
 H=
 \begin{bmatrix}
 h&h&\cdots&h\\
 h^2(K-\tfrac12)&h^2(K-\tfrac32)&\cdots&h^2/2
 \end{bmatrix},\qquad
 \mathbf b_i=\begin{bmatrix}v_i^f-v_i^0\\g_i-s_i-T_f v_i^0\end{bmatrix}.
 \label{eq:moments}
\end{equation}
The minimum Euclidean-norm correction satisfying $H\bar{\mathbf a}_i=\mathbf b_i$ is
\begin{equation}
 \bar{\mathbf a}_i=\mathbf a_i+
 H^\top(HH^\top)^{-1}(\mathbf b_i-H\mathbf a_i).
 \label{eq:correction}
\end{equation}
The two rows of $H$ are independent for $K\geq2$, so only a $2\times2$ linear system is required for each coordinate. The correction requires no additional iterative optimization solve.

Starting from the prescribed initial state, the corrected controls generate
\begin{subequations}\label{eq:rollout}\begin{align}
 \bar\bv_{i,k+1}&=\bar\bv_{i,k}+h\bar\ba_{i,k},\\
 \bar\bp_{i,k+1}&=\bar\bp_{i,k}+h\bar\bv_{i,k}+\tfrac12h^2\bar\ba_{i,k}.
\end{align}\end{subequations}
For $0\leq\tau\leq h$, the returned interval motion is
\begin{equation}
 \bar\bp_i(t_k+\tau)=\bar\bp_{i,k}+\tau\bar\bv_{i,k}
 +\tfrac12\tau^2\bar\ba_{i,k},\qquad
 \bar\bv_i(t_k+\tau)=\bar\bv_{i,k}+\tau\bar\ba_{i,k}.
 \label{eq:arc}
\end{equation}
It follows immediately from~\eqref{eq:rollout} that position and velocity agree at neighboring interval boundaries. Equation~\eqref{eq:correction} enforces both terminal moments, giving the prescribed final position and zero final velocity. The output therefore satisfies~\eqref{eq:continuous} on every interval, with continuous velocity and bounded, possibly discontinuous acceleration when the subsequent limit checks pass.

To relate reconstruction to numerical recovery, define
\begin{equation}
 D_i=\sum_{k=0}^{K-1}\norm{\mathbf E_{i,k}}_2^2,\qquad
 \mathbf e_i=\sum_{k=0}^{K-1}\mathbf E_{i,k},\qquad
 \norm{\mathbf e_i}_2\leq\sqrt{K D_i}.
 \label{eq:terminal-bound}
\end{equation}
The inequality follows from Cauchy--Schwarz. With exact candidate endpoint bounds and $\ba_{i,k}=(\bv_{i,k+1}-\bv_{i,k})/h$, the uncorrected terminal velocity is exact, and the remaining displacement mismatch is $\mathbf e_i$. For any coordinate $d$,~\eqref{eq:correction} therefore gives
\begin{equation}
 \norm{\bar{\mathbf a}_{i,d}-\mathbf a_{i,d}}_2
 \leq\left\|H^\top(HH^\top)^{-1}
 \begin{bmatrix}0\\1\end{bmatrix}\right\|_2\sqrt{K D_i}.
 \label{eq:correction-bound}
\end{equation}
For a fixed grid, the bound vanishes with $D_i$; the actual correction also depends on cancellation among interval defects. The integrated motion is checked against every obstacle and UAV pair, together with the speed, acceleration, and workspace bounds. These checks account for its displacement from the QP candidate while retaining the prescribed horizon and endpoint states.

\begin{algorithm2e}[H]
\caption{LIOM for fixed-horizon cooperative trajectory optimization}
\label{alg:liom}
\KwIn{Starts and goals, obstacle map, horizon $T_f$, and motion limits}
\KwOut{An accepted cooperative trajectory or planning failure}
Generate synchronized reference $\widehat{\mathbf P},\widehat{\mathbf V}$\;
\If{reference generation fails}{\Return{planning failure}\;}
Initialize $\rho_q\gets\rho_q^{(0)}$ and $\rho_s\gets\rho_s^{(0)}$\;
\For{$\nu=0,\ldots,\nu_{\max}-1$}{
  \If{the computation budget is exhausted}{\textbf{break}\;}
  Construct spatial boxes $\Gamma_{i,k}$ and the active set $\mathcal A$\;
  \If{spatial boxes are invalid or exclude a fixed endpoint}{\textbf{break}\;}
  Construct relative boxes $\mathcal R_{ij,k}$ for $(i,j,k)\in\mathcal A$\;
  \If{the nominal computation budget is exhausted}{\textbf{break}\;}
  Solve~\eqref{eq:subproblem} for $\mathbf P,\mathbf V,\boldsymbol\delta$\;
  \If{no finite candidate is available or a solver exception occurs}{\textbf{break}\;}
  Evaluate $D$, $S$, and the hard-constraint violation\;
  \If{$D\leq\varepsilon_q$, $S\leq\varepsilon_s$, and the hard constraints hold}{
    Correct endpoint control moments using~\eqref{eq:correction}\;
    Integrate the corrected controls using~\eqref{eq:rollout}\;
    \If{the reconstructed trajectory passes complete-interval clearance and motion checks}{
      \Return{$(\bar{\mathbf P},\bar{\mathbf V},\bar{\mathbf A})$}\;
    }
  }
  Update $\rho_q,\rho_s$ independently using~\eqref{eq:weights}\;
  Set $\widehat{\mathbf P}\gets\mathbf P$, $\widehat{\mathbf V}\gets\mathbf V$\;
}
\Return{planning failure}\;
\end{algorithm2e}

\subsection{Computational Structure}
Acceleration elimination leaves $6N(K+1)$ state components, including the fixed boundary components. Relative-corridor slack introduces at most $C_N^2(K+1)$ auxiliary variables, where $C_N^2=N(N-1)/2$ is the number of UAV pairs. Each active pair--node triple contributes six one-sided linear inequalities. Geometric screening removes unnecessary pair restrictions, with worst-case relative-constraint growth of $O(N^2K)$.

Within each solve, the constraint Jacobian and objective Hessian are constant and sparse, reflecting temporal and pairwise coupling. Warm starts transfer information across successive corridor reconstructions.
